\documentclass[12pt]{article}
\usepackage[utf8]{inputenc}
\usepackage[T1]{fontenc}
\usepackage{hyperref}
\title{Yangben Language Resources for K-12 Curriculum Development}
\author{Compiled by Claude Sonnet 5.5 (claude-sonnet-5-5) for Andr\'as Kornai}
\date{}
\begin{document}
\maketitle

\section{Introduction}
Yangben (ISO 639-3: yav; Glottocode yang1293) is an Atlantic-Congo (Bantu) language of the Centre region of Cameroon. Speakers live mainly in Yangben canton, south of Bokito. Global Recordings Network gives about 2,300 speakers. Other sources give higher figures. Alternate names include Nuasue, Nukalonge, Kalong, Yambassa and Central Yambassa. Webonary lists three dialects: Nukalonge, Nutanga and Numende. The searches found few materials. Nothing found is a school textbook or primer.

\section{Primary Educational Resources}
No primers, textbooks or school curricula in Yangben turned up in the searches. The reference description says the language is taught as a subject in some schools. No programme documentation was found to confirm this.

\section{Linguistic and Reference Materials}
\begin{itemize}
\item Yangben--French dictionary (SIL Cameroon, 1,891 entries, 2016, described as a rough draft): \url{https://www.silcam.org/resources/archives/84830}
\item Webonary online Yangben dictionary, with an introduction covering dialects, classification and language use: \url{https://www.webonary.org/yangben/overview/introduction/?lang=en}
\end{itemize}

\section{Religious and Community Resources}
\begin{itemize}
\item Scripture Earth listing for Yangben (Nuasue), including a New Testament: \url{https://www.scriptureearth.org/00eng.php?iso=yav}
\item Joshua Project people-group profile: \url{https://joshuaproject.net/people_groups/15979}
\end{itemize}

\section{Digital and Audio Resources}
\begin{itemize}
\item Global Recordings Network page for Yangben, which lists audio materials: \url{https://globalrecordings.net/en/language/yav}
\end{itemize}
Warning: Localazy pages for ``yav yangben'' describe a Chinese language and should not be used for this language.

\section{Curriculum Development Context}
The language is endangered. Adults use it, but not all young people do, and French is widespread. A curriculum would have to start from the draft dictionary and the New Testament, with new reading materials written with community speakers. Little exists in the sources found, so developers should treat the materials above as a starting point and consult speakers and SIL Cameroon.

\end{document}